\documentclass[12pt]{article}

\usepackage[T1]{fontenc}
\usepackage{newpxtext,newpxmath}
\usepackage{amsmath}
\usepackage{geometry}
\usepackage{microtype}
\usepackage{enumitem}
\usepackage{xcolor}
\usepackage{fancyhdr}
\usepackage{comment}
\usepackage[most]{tcolorbox}

\newif\ifsolutions
\solutionstrue   % flip to \solutionsfalse for student version
%   \solutionsfalse

\tcbset{
  solutionstyle/.style={
    enhanced,
    breakable,
    frame hidden,
    borderline west={1pt}{0pt}{blue!60!black},
    colback=white,
    sharp corners,
    left=10pt,
    right=0pt,
    top=6pt,
    bottom=6pt,
    parskip=0pt,
    fontupper=\setlength{\parindent}{0pt}\sffamily,
    before upper=\noindent,
  }
}

\newenvironment{solution}{
  \ifsolutions
     \begin{tcolorbox}[solutionstyle]
     \textbf{Solution.}\\[0.3em]
  \else
     \comment
  \fi
}{
  \ifsolutions
     \end{tcolorbox}
  \else
     \endcomment
  \fi
}

\geometry{margin=1in}

% ----- Header & footer -----
\pagestyle{fancy}
\fancyhf{}
\lhead{\textbf{ECON 1001}}
\chead{Problem Set 2 -- Part B}
\rhead{due Jan.\ 23, 2026, 6\,pm}
\renewcommand{\headrulewidth}{0.4pt}
\cfoot{\thepage}
\renewcommand{\footrulewidth}{0pt}
\setlength{\headheight}{10.5pt} % avoid fancyhdr warning
\addtolength{\topmargin}{-4pt}

% \fancypagestyle{firstpage}{
%   \fancyhf{}
%   \lhead{\textbf{ECON 4433}}
%   \chead{Problem Set 3 (v2)}
%   \rhead{due 9 Dec.\ 2025}
%   \renewcommand{\headrulewidth}{0.4pt}
%   \cfoot{\thepage}
% }

\setlength{\parskip}{0.6em}
\setlength{\parindent}{0pt}

% Enumerate styles if needed:
\setlist[enumerate,1]{label=\textbf{\arabic*.}}
\setlist[enumerate,2]{label=\textbf{(\alph*)}}
% Reduce spacing in enumerates
\setlist[enumerate]{itemsep=0.3em, topsep=0.3em}
\setlist[itemize]{itemsep=0.2em, topsep=0.2em}

\begin{document}
% \thispagestyle{firstpage}
\noindent\textbf{Instructions.} Problem Set 2 has two parts: Part A (Achieve) and Part B (this problem). Part B counts for 10 percent of your total Problem Set 2 score.

You may work in a group, but each student must submit their own explanation via Canvas.

The goal of this problem is to think carefully about timing, saving, and opportunity cost.

To answer this problem carefully, you will need to compare the value of
different streams of payments over time. You should perform whatever
calculations are necessary to support your conclusions; however, your grade will be based primarily on the clarity and correctness of your economic reasoning rather than on exact numerical precision. \textbf{This problem will be graded with a light touch: a thoughtful, well-explained approach earns full credit even if your calculations are not exact. Short explanations (two or three sentences) are sufficient.}

If you make additional assumptions, state them clearly and proceed.
Reasonable assumptions will not be penalized.

You may find it helpful to use a spreadsheet to compute and organize future
values of savings under each option.

After grading, we may ask permission to share one or more especially clear
student submissions as an informal answer key and study guide.

\noindent\textbf{Submission.}
Upload a single PDF containing your written explanation to Canvas under the
\emph{Problem Set 2 -- Part B} assignment. Typed or neatly handwritten submissions are both acceptable. If you wish, you may also upload a spreadsheet showing your calculations.

\begin{enumerate}

    \item \textbf{Mortgage Choice and Saving.}

          It is a few years down the road, and you have settled into a stable career. For the rest of your life, after covering all non-housing expenses, you will have \$15{,}000 per year available to spend on housing. You decide to buy your dream home, which requires borrowing money from a bank. The bank offers you two mortgage options that allow you to purchase the home today:

          \begin{itemize}
              \item \textbf{The long mortgage:} You pay \$10{,}000 each year for 30 years.
              \item \textbf{The short mortgage:} You pay \$15{,}000 each year for 15 years.
          \end{itemize}

          Any portion of your \$15{,}000 annual housing budget that you do not spend on mortgage
          payments is deposited into a savings account. Your savings account earns a constant
          annual return of 7 percent. For simplicity, assume there is no inflation and your
          income does not change over time.

          \begin{enumerate}
              \item After 30 years, what is the total amount of money you will have paid to the bank
                    under each mortgage option?

              \item Which mortgage do you choose? Briefly explain your reasoning.

              \item How would your answer change if the savings account instead earned a return of
                    5 percent per year? Briefly explain.
          \end{enumerate}

\end{enumerate}
\end{document}